\subsection{Experimental Results}
\label{sec:results}
% \subsection{Results}

%TODO: WHY WORD2VEC is worst

Table~\ref{tab:type_iden_perf} (\textit{Original Performance} tab) shows the performance of the studied VTI approaches. 
Note that all our experiments were conducted on a workstation with a P100 GPU, dual vCPUs, and 32GB RAM.


% % Please add the following required packages to your document preamble:
% % \usepackage{multirow}
% \begin{table}[]
% \centering
% \caption{Vulnerability type identification performance}
% \label{tab:type_iden_perf}
% \begin{tabular}{ll|r|r|r}
% \toprule
% \multicolumn{2}{l|}{\textit{Metrics}}                                           & \base    & Word2vec & CodeBERT \\ \midrule
% \multicolumn{1}{l|}{\multirow{3}{*}{\textit{Micro-avg}}}    & \textit{Prec.}    & 0.71     & 0.55         & 0.73     \\  
% \multicolumn{1}{l|}{}                              & \textit{Rec.}              & 0.69     & 0.43         & 0.73     \\ 
% \multicolumn{1}{l|}{}                              & \textit{F1}                & 0.70     & 0.48         & 0.73     \\ \midrule
% \multicolumn{1}{l|}{\multirow{3}{*}{\textit{Macro-avg}}}    & \textit{Prec.}    & 0.65     & 0.41         & 0.67     \\  
% \multicolumn{1}{l|}{}                              & \textit{Rec.}              & 0.60     & 0.24         & 0.65     \\ 
% \multicolumn{1}{l|}{}                              & \textit{F1}                & 0.62     & 0.28         & 0.65     \\ \midrule
% \multicolumn{1}{l|}{\multirow{3}{*}{\textit{Weighted-avg}}} & \textit{Prec.}    & 0.71     & 0.50         & 0.73     \\ 
% \multicolumn{1}{l|}{}                              & \textit{Rec.}              & 0.69     & 0.43         & 0.73     \\ 
% \multicolumn{1}{l|}{}                              & \textit{F1}                & 0.70     & 0.43         & 0.73     \\ \midrule
% \multicolumn{1}{l|}{\multirow{3}{*}{\textit{Sample-avg}}}   & \textit{Prec.}    & 0.65     & 0.47         & 0.71     \\ 
% \multicolumn{1}{l|}{}                              & \textit{Rec.}              & 0.68     & 0.42         & 0.72     \\ 
% \multicolumn{1}{l|}{}                              & \textit{F1}                & 0.65     & 0.42         & 0.70     \\ \bottomrule \toprule
% \multicolumn{2}{l|}{\textit{ExactMatchRatio}}                                   & 0.54     & 0.30         & 0.57     \\ \midrule
% \multicolumn{2}{l|}{\textit{Hamming Score}}                                     & 0.62     & 0.39         & 0.66     \\ \midrule
% \multicolumn{2}{l|}{\textit{Accuracy}}                                          & 0.89     & 0.84         & 0.90     \\ \midrule
% \multicolumn{2}{l|}{\textit{Training time (min)}}                               & 2.48     & 43.5         & 82.3     \\ \midrule
% \multicolumn{2}{l|}{\textit{Predict. time (sec)}}                               & 4.0     & 12.3         & 22.3     \\ \midrule
% \multicolumn{2}{l|}{\textit{GPU required}}                                      & No     & Yes           & Yes     \\
% \bottomrule
% \end{tabular}
% \end{table}






\begin{table*}[]
\centering
\caption{Vulnerability type identification performance}
\label{tab:type_iden_perf}
\begin{tabular}{@{}ll|rrr|rrr@{}}
\toprule
\multicolumn{2}{l}{} & \multicolumn{3}{c|}{\textit{Original Performance}}      & \multicolumn{3}{c}{\textit{After Enhancement}}       \\ \cmidrule(l){3-8} 
\multicolumn{2}{l}{}                         & \base    & Word2vec  & CodeBERT  & \base    & Word2vec  & CodeBERT  \\ \cmidrule(l){1-8} 
\multirow{3}{*}{\textit{micro-avg}}        & Prec.    & 0.71    & 0.55      & 0.73      & 0.72    & 0.60      & 0.74      \\
                                  & Rec.     & 0.69    & 0.43      & 0.73      & 0.78    & 0.56      & 0.79      \\
                                  & F1       & 0.70    & 0.48      & 0.73      & 0.75    & 0.58      & 0.77      \\
                                  \midrule
\multirow{3}{*}{\textit{macro-avg}}        & Prec.    & 0.65    & 0.41      & 0.67      & 0.69    & 0.70      & 0.70      \\
                                  & Rec.     & 0.60    & 0.24      & 0.65      & 0.75    & 0.48      & 0.76      \\
                                  & F1       & 0.62    & 0.28      & 0.65      & 0.71    & 0.54      & 0.72      \\
                                  \midrule
\multirow{3}{*}{\textit{weighted-avg}}     & Prec.    & 0.71    & 0.50      & 0.73      & 0.73    & 0.61      & 0.74      \\
                                  & Rec.     & 0.69    & 0.43      & 0.73      & 0.78    & 0.56      & 0.79      \\
                                  & F1       & 0.70    & 0.47      & 0.73      & 0.75    & 0.55      & 0.76      \\
                                  \midrule
\multirow{3}{*}{\textit{sample-avg}}       & Prec.    & 0.65    & 0.47      & 0.71      & 0.73    & 0.57      & 0.75      \\
                                  & Rec.     & 0.68    & 0.42      & 0.72      & 0.80    & 0.60      & 0.81      \\
                                  & F1       & 0.65    & 0.42      & 0.70      & 0.74    & 0.55      & 0.76      \\
                                  \midrule
\multicolumn{2}{l|}{\textit{ExactMatchRatio}}          & 0.54    & 0.30      & 0.57      & 0.60    & 0.32      & 0.57      \\
\multicolumn{2}{l|}{\textit{HammingScore}}            & 0.62    & 0.39      & 0.66      & 0.69    & 0.49      & 0.71      \\
\multicolumn{2}{l|}{\textit{Accuracy}}                 & 0.89    & 0.84      & 0.90      & 0.91    & 0.86      & 0.91      \\
\midrule
\multicolumn{2}{l|}{\textit{Training time (ms)}}            & 148,824 & 2,610,032 & 4,938,020 & 148,851 & 2,610,035 & 4,938,035 \\
\multicolumn{2}{l|}{\textit{Predict. time (ms)}}            & 4,320   & 12,349    & 22,282    & 5,243   & 13,841    & 24,028    \\
\midrule
\multicolumn{2}{l|}{\textit{GPU required}}             & No      & Yes       & Yes       & No      & Yes       & Yes       \\ \bottomrule
\end{tabular}
\end{table*} 

As seen in \textit{Original Performance} tab of Table~\ref{tab:type_iden_perf}, VTI with CodeBERT achieved the best performance among the three approaches. 
%
%
However, the performance of this deep pre-trained model is slightly better than that of \base, about 4\%-7.5\% in F1-scores for \textit{macro average}, \textit{micro average}, \textit{weight average}, and \textit{sample average}. Meanwhile, this advanced method improves 5.5\%, 6.4\%, and only 1.1\% in \textit{exact match ratio}, \textit{hamming score}, and \textit{accuracy}, respectively. \base even significantly outperforms the VTI approach with Word2vec in all the considering metrics. Especially, the micro-average of \base doubles the corresponding figure of Word2vec. 
%
%



Table~\ref{tab:perf_type} shows the detailed classification performance (F1-score) of the three approaches for each vulnerability type. \base achieves F1-scores comparable with those of the CodeBERT-based approach, even slightly better for \textit{Memory Corruption} and \textit{Others}. Additionally, \base significantly improves the performance of the Word2vec-based method for all types. 
%
%
These results empirically demonstrate that \textit{\base with a traditional multi-label classification model and the old-fashioned TF-IDF could achieve a very competitive performance compared to the advanced VTI approaches}.

Especially, \base is much more efficient than the other approaches. Specially, \base is trained \textbf{20X} and \textbf{40X faster} than the VTI approaches based on Word2vec and CodeBERT. Moreover, \base is also much faster in predicting vulnerability type of vulnerable functions, 3X and 5X faster compared to the Word2vec-based and CodeBERT-based models. Notably, unlike the neural network approaches, \base can be trained and tested \textbf{without GPU}, while the Word2vec-based and CodeBERT-based models require GPU and consume much computational resource. This advantage enables a flexible deployment for \base on developers' machines which might not be very powerful with limited resources while retaining a competitive VIT performance.



% \begin{table}[]
% \centering
% \caption{Performance in F1 score by types}
% \label{tab:perf_type}
% \begin{tabular}{l|r|r|r}
% \toprule
% \textit{Types}  & \base     & Word2vec & CodeBERT \\ \midrule
% +Info           & 0.41      & 0.00         & \textbf{0.47}         \\ \midrule
% +Priv           & 0.63      & 0.13         & \textbf{0.67}         \\ \midrule
% Bypass          & 0.42      & 0.00         & \textbf{0.48}         \\ \midrule
% DoS             & 0.81      & 0.68         & \textbf{0.83}         \\ \midrule
% Exec. Code      & 0.64      & 0.34         & \textbf{0.66}         \\ \midrule
% Mem. Corr.      & \textbf{0.70}      & 0.36         & 0.68         \\ \midrule
% Overflow        & 0.67      & 0.37         & \textbf{0.75}         \\ \midrule
% Others          & \textbf{0.69}      & 0.34         & 0.68         \\ \bottomrule
% \end{tabular}
% \end{table}


\begin{table}[]
\centering
\caption{Performance in F1 score by types}
\label{tab:perf_type}
\begin{tabular}{l|r|r|r}
\toprule
\textit{Types}  & \base     & Word2vec & CodeBERT \\ \midrule
Information Gain           & 0.41      & 0.00         & \textbf{0.47}         \\ \midrule
Privilege Gain           & 0.63      & 0.13         & \textbf{0.67}         \\ \midrule
Bypass          & 0.42      & 0.00         & \textbf{0.48}         \\ \midrule
DoS             & 0.81      & 0.68         & \textbf{0.83}         \\ \midrule
Execution Code      & 0.64      & 0.34         & \textbf{0.66}         \\ \midrule
Memory Corruption      & \textbf{0.70}      & 0.36         & 0.68         \\ \midrule
Overflow        & 0.67      & 0.37         & \textbf{0.75}         \\ \midrule
Others          & \textbf{0.69}      & 0.34         & 0.68         \\ \bottomrule
\end{tabular}
\end{table}

% \subsection{Discussion}


